% \documentclass{beamer}
%Для защит онлайн лучше использовать разрешение 16x9
\documentclass[aspectratio=169]{beamer}

%%% PDF settings
\pdfvariable minorversion 7 % Set PDF version to 1.7.

%%% Fonts and language setup.
\usepackage{polyglossia}
% Setup fonts.
\usepackage{fontspec}
\setmainfont{CMU Serif}
\setsansfont{CMU Sans Serif}
\setmonofont{CMU Typewriter Text}

\usepackage{microtype} % Add fancy-schmancy font tricks

\usepackage{xcolor} % Add colors support.

%% Math
\usepackage{amsmath, amsfonts, amssymb, amsthm, mathtools} % Advanced math tools.
\usepackage{thmtools}
\usepackage{unicode-math} % Allow TTF and OTF fonts in math and allow direct typing unicode math characters.
\unimathsetup{
    warnings-off={
            mathtools-colon,
            mathtools-overbracket
        }
}
\setmathfont{Latin Modern Math} % default
\setmathfont[range={\setminus,\varnothing,\smashtimes}]{Asana Math}

%%% Images
\usepackage{graphicx}
\graphicspath{{figures/}}
\usepackage{import}

%%% Polyglossia setup after (nearly) everything as described in documentation.
\setdefaultlanguage{russian}
\setotherlanguage{english}

\usepackage{csquotes}

%%% Custom commands
\newcommand{\R}{\mathbb{R}}
\newcommand{\N}{\mathbb{N}}
\newcommand{\Z}{\mathbb{Z}}
\newcommand{\Q}{\mathbb{Q}}
\newcommand{\C}{\mathbb{C}}
\newcommand{\id}{\mathrm{id}}
\AtBeginDocument{\renewcommand{\leq}{\leqslant}}
\AtBeginDocument{\renewcommand{\geq}{\geqslant}}
\AtBeginDocument{\renewcommand{\Re}{\operatorname{Re}}}
\AtBeginDocument{\renewcommand{\Im}{\operatorname{Im}}}
\AtBeginDocument{\renewcommand{\phi}{\varphi}}
\AtBeginDocument{\renewcommand{\epsilon}{\varepsilon}}

%%% theorem-like envs
\theoremstyle{definition}

\declaretheoremstyle[spaceabove=0.5\topsep,
    spacebelow=0.5\topsep,
    headfont=\bfseries\sffamily,
    bodyfont=\normalfont,
    headpunct=.,
    postheadspace=5pt plus 1pt minus 1pt]{myStyle}
\declaretheoremstyle[spacebelow=\topsep,
    headfont=\bfseries\sffamily,
    bodyfont=\normalfont,
    headpunct=.,
    postheadspace=5pt plus 1pt minus 1pt,]{myStyleWithFrame}
\declaretheoremstyle[spacebelow=\topsep,
    headfont=\bfseries\sffamily,
    bodyfont=\normalfont,
    headpunct=.,
    postheadspace=5pt plus 1pt minus 1pt,
    qed=\blacksquare]{myProofStyleWithFrame}

\usepackage[breakable]{tcolorbox}
\tcbset{sharp corners=all, colback=white}
% \tcolorboxenvironment{theorem}{}
% \tcolorboxenvironment{theorem*}{}
% \tcolorboxenvironment{axiom}{}
% \tcolorboxenvironment{assertion}{}
% \tcolorboxenvironment{lemma}{}
% \tcolorboxenvironment{proposition}{}
% \tcolorboxenvironment{corollary}{}
% \tcolorboxenvironment{definition}{}
% \tcolorboxenvironment{proofReplace}{toprule=0mm,bottomrule=0mm,rightrule=0mm, colback=white, breakable }

\declaretheorem[name=Теорема, style=myStyleWithFrame]{theorem}
\declaretheorem[name=Теорема, numbered=no, style=myStyleWithFrame]{theorem*}
\declaretheorem[name=Аксиома, sibling=theorem, style=myStyleWithFrame]{axiom}
\declaretheorem[name=Преположение, sibling=theorem, style=myStyleWithFrame]{assertion}
\declaretheorem[name=Лемма, style=myStyleWithFrame]{lemma}
\declaretheorem[name=Предложение, sibling=theorem, style=myStyleWithFrame]{proposition}
\declaretheorem[name=Следствие, numberwithin=theorem, style=myStyleWithFrame]{corollary}

\declaretheorem[name=Определение, style=myStyleWithFrame]{definition}
\declaretheorem[name=Свойство, style=myStyle]{property}
\declaretheorem[name=Свойства, numbered=no, style=myStyle]{propertylist}

\declaretheorem[name=Пример, style=myStyle]{example}
\declaretheorem[name=Замечание, numbered=no, style=myStyle]{remark}

\declaretheorem[name=Доказательство, numbered=no, style=myProofStyleWithFrame]{proofReplace}
\renewenvironment{proof}[1][\proofname]{\begin{proofReplace}}{\end{proofReplace}}
% \declaretheorem[name=Доказательство, numbered=no, style=myProofStyleWithFrame]{longProof}

%%% Memoir settings
\chapterstyle{ger}
\setlength{\headheight}{2\baselineskip}

%%% HyperRef
\usepackage{hyperref}


% То, что в квадратных скобках, отображается внизу по центру каждого слайда.
\title[Обзор инфраструктуры сборки]{Обзор инфраструктуры сборки проектов с открытым исходным кодом}

% То, что в квадратных скобках, отображается в левом нижнем углу.
\institute[СПбГУ]{}

% То, что в квадратных скобках, отображается в левом нижнем углу.
\author[Николай Пономарев]{Пономарев Николай Алексеевич, группа 21.Б10-мм}

\begin{document}
{
\setbeamertemplate{footline}{}
% Лого университета или организации, отображается в шапке титульного листа
\begin{frame}
	\includegraphics[width=1.4cm]{pictures/SPbGU_Logo.png}
	\vspace{-35pt}
	\hspace{-10pt}
	\begin{center}
		\begin{tabular}{c}
			\scriptsize{Санкт-Петербургский государственный университет} \\
			\scriptsize{Кафедра системного программирования}
		\end{tabular}
		\titlepage
	\end{center}

	\btVFill

	{\scriptsize
		% У научного руководителя должна быть указана научная степень
		\textbf{Научный руководитель:} Я. А. Кириленко, старший преподаватель кафедры системного программирования
	}
	\begin{center}
		\vspace{5pt}
		\scriptsize{Санкт-Петербург\\
			2024}
	\end{center}

\end{frame}
}

\begin{frame}[fragile]
	\frametitle{Введение}
	\begin{itemize}
		\item На рынке всё больше устройств на архитектурах отличных от \amd{}, например \riscv{} или \arm{}
		\item Нужно переносить используемые приложения, но сначала надо научиться собирать приложения на другой архитектуре и настраивать \ci{} для них
		\item В основном изучаем проекты на \textsc{C/C++}
	\end{itemize}
\end{frame}

\begin{frame}
	\frametitle{Постановка задачи}
	\textbf{Целью} данной работы является изучение возможностей по созданию
	инфраструктуры для сборки проектов с открытым исходным кодом.

	\textbf{Задачи}:
	\begin{enumerate}
		\item Провести обзор видов сборки проектов на языках \textsc{C}/\textsc{C++}
		\item Провести обзор поддержки нативной и кросс-сборки в системах сборки
		\item Исследовать возможность настройки \ci{} на архитектурах, отличных от \amd{}
		\item Протестировать изученные возможности систем сборки на реальных проектах
		\item Описать процесс компиляции и настройки \ci{} на примере проектов с различными системами сборки
	\end{enumerate}

\end{frame}

\begin{frame}
	\frametitle{Виды сборки}

	\begin{itemize}
		\item Нативная сборка
		      \begin{itemize}
			      \item Сборка производится на целевой платформе
			      \item Обычный способ сборки
			      \item Требует оборудования на целевой архитектуре
		      \end{itemize}
		\item Кросс-сборка
		      \begin{itemize}
			      \item Целевая платформа отличается от платформы, на которой производится сборка
			      \item Могут быть сложности с удовлетворением зависимостей
			      \item Не во всех проектах есть поддержка
		      \end{itemize}
		\item Нативная сборка в эмуляторе
		      \begin{itemize}
			      \item Приложения считают, что они запущены на своей целевой платформе
			      \item \qemu{} + \textsc{binfmt\_misc} позволяют не запускать виртуальную машину целиком
			      \item Малая скорость исполнения
		      \end{itemize}
	\end{itemize}
\end{frame}

\begin{frame}
	\frametitle{Поддержка кросс-компиляции в системах сборки}

	\begin{itemize}
		\item Make
		      \begin{itemize}
			      \item При аккуратно написанном Makefile поддерживает кросс-компиляцию
		      \end{itemize}
		\item Autotools
		      \begin{itemize}
			      \item Поддерживает кросс-компиляцию из коробки
			      \item Могут быть сложности с указанием зависимостей
		      \end{itemize}
		\item CMake
		      \begin{itemize}
			      \item Хорошая поддержка кросс-сборки
			      \item Иногда требует правки сборочных скриптов
		      \end{itemize}
		\item QMake
		      \begin{itemize}
			      \item Требует сборки в специальном кросс режиме
			      \item Можно использовать нативную сборку в эмуляторе
		      \end{itemize}
		\item Meson
		      \begin{itemize}
			      \item Большие возможности кросс-компиляции
			      \item Плохая документация
			      \item Сложности с поиском зависимостей с помощью pkg-config
		      \end{itemize}
	\end{itemize}
\end{frame}

\begin{frame}
	\frametitle{Поддержка в \ci{} системах}

	\begin{itemize}
		\item \GitHub{}
		      \begin{itemize}
			      \item В своём облаке поддерживает только \amd{}
			      \item Self-hosted runner официально поддерживает \amd{} и \arm{}
			      \item Self-hosted runner написан на \dotnet{}, проверить работу на \riscv{} невозможно
		      \end{itemize}
		\item \gitlab{}
		      \begin{itemize}
			      \item В своём облаке поддерживает только \amd{}
			      \item Self-hosted runner поддерживает множество архитектур, в том числе \riscv{}
		      \end{itemize}
		\item \jenkins{}
		      \begin{itemize}
			      \item Написан на \java{}, которая поддерживает множество архитектур
		      \end{itemize}
	\end{itemize}
\end{frame}

\begin{frame}
	\frametitle{Тестирование}

	\begin{itemize}
		% \item Libusb (Autotools)
		\item SingleApplication (CMake)
		      \begin{itemize}
			      \item \cmake{} пытается использовать утилиты \qt{} для целевой платформы
		      \end{itemize}
		      % \item TRIK Runtime (QMake)
		      % \item Umockdev (Meson)
		      % \begin{itemize}
		      % 	\item
		      % \end{itemize}
		\item TRIK Runtime (Jenkins)
		      \begin{itemize}
			      \item Попытка запустить на StarFive VisionFive V2
			      \item Потребовалась пересборка ядра
			      \item На одной плате серверу и агенту \enquote{тесно}
		      \end{itemize}
	\end{itemize}

\end{frame}

\begin{frame}
	\frametitle{Результаты}
	Полностью достигнуты следующие результаты:
	\begin{enumerate}
		\item Проведен обзор видов сборки проектов на языках \textsc{C}/\textsc{C++}
		\item Проведен обзор поддержки нативной и кросс-сборки в системах сборки
		\item Исследована возможность настройки \ci{} на архитектурах, отличных от \amd{}
		\item Протестированы изученные возможности систем сборки на реальных проектах
	\end{enumerate}
	Частично достигнуты следующие результаты:
	\begin{enumerate}
		\item Описан процесс компиляции и настройки \ci{} на примере проектов с различными системами сборки
	\end{enumerate}

\end{frame}
\end{document}
